\documentclass[11pt]{article}
\usepackage[T1]{fontenc}
\usepackage{amsmath,amssymb,amsthm,booktabs,geometry,hyperref}
\geometry{margin=1in}
\emergencystretch=2em
\hypersetup{pdftitle={The seventh ternary density Hales--Jewett number},pdfauthor={P. Graham},pdfsubject={Computer-assisted proof of the exact value 1302},colorlinks=true,urlcolor=blue,citecolor=blue,linkcolor=blue}
\newtheorem{theorem}{Theorem}
\newtheorem{lemma}{Lemma}
\newcommand{\cube}[1]{\{0,1,2\}^{#1}}
\title{The seventh ternary density Hales--Jewett number}
\author{P. Graham\\ \small\url{https://orcid.org/0009-0007-8493-1505}}
\date{October 20, 2026}
\begin{document}
\maketitle
\begin{abstract}
For a word $w\in\{0,1,2,*\}^n$ with at least one wildcard, its combinatorial line is $\{w(0),w(1),w(2)\}$, where $w(a)$ replaces every wildcard by $a$. Let $c_{n,3}$ be the maximum size of a subset of $\{0,1,2\}^n$ containing no combinatorial line. We give a computer-assisted proof that $c_{7,3}=1302$, closing the gap in Polymath1's bounds $1302\leq c_{7,3}\leq1348$, recorded by Tao in 2009 and O'Bryant in 2015, and confirming Polymath's tentative conjecture. The supplied witness is a realization of Polymath's known lower bound. Lower-dimensional bounds and a global alphabet permutation reduce the upper bound to nineteen propositional unsatisfiability obligations: three auxiliary bounds and sixteen slice cases. A separately written semantic auditor checks every clause under one simultaneous counting interpretation. All nineteen LRAT refutations were checked with FRAT-rs and drat-trim's LRAT checker; eleven also with the formally verified checker cake\_lpr. The reduction from sets to formulas is not formalized.
\end{abstract}

The proofs, the search code, the certificate checkers and the text were produced by AI models (from OpenAI and Anthropic), working under the author's direction.
The result has not been reviewed by human experts.

\section{Introduction}
The density Hales--Jewett theorem of Furstenberg and Katznelson~\cite{fk}, with a later combinatorial proof by Polymath~\cite{annals}, concerns large subsets of word spaces and the combinatorial lines they must contain. A combinatorial line varies a nonempty set of coordinates through the same letter while fixing the rest. The finite extremal values are not determined by the asymptotic theorem. For a fixed dimension $n$, write $c_{n,3}$ for the maximum size of a line-free subset of $\cube n$. D.H.J. Polymath determined these numbers through dimension six and established the lower bound 1302 in dimension seven. Section 2 of its 2010 paper calls this the bound ``which we tentatively conjecture to be the correct bound''~\cite{polymath}. OEIS A156989, accessed October 20, 2026, lists the established terms through dimension six, ending at 450~\cite{oeis}.

The historical gap between Polymath1's bounds is explicitly documented. In his March 4, 2009 Polymath post, Tao writes: ``The current best-known bounds for $c_7$ are $1302\leq c_7\leq1348$.''~\cite{tao2009} O'Bryant's 2015 paper repeats ``$1302\leq c_{7,3}\leq1348$'' in its list of density Hales--Jewett numbers~\cite[Section 3.3, p.~7]{obryant}. We close this gap by proving the matching unrestricted upper bound. To our knowledge, this is the first proof of $c_{7,3}\leq1302$. Searches of the primary literature, recent citing works, OEIS and current postings through October 20, 2026 found no earlier proof of this upper bound.

\begin{theorem}\label{main}
Every subset of $\cube 7$ with no combinatorial line has at most 1302 points, and a line-free subset of size 1302 exists. Equivalently, $c_{7,3}=1302$.
\end{theorem}

The contribution claimed here is the upper bound for arbitrary subsets. The known lower bound and the conjectured value belong to Polymath. The global alphabet permutation only orders slice labels without loss of generality; the set need not be invariant under it. The argument proves the caps $c_1\leq2$, $c_4\leq52$ and $c_5\leq150$ by auxiliary SAT refutations, and obtains $c_2\leq6$, $c_3\leq18$ and $c_6\leq450$ by slicing. Ordering the three first-coordinate slice sizes leaves sixteen cases, indexed by the largest size from 435 through 450.

Historically, all nineteen CaDiCaL-generated text LRAT refutations were checked by FRAT-rs's \texttt{lratchk} and drat-trim's standalone \texttt{lrat-check}; cake\_lpr also checked the three auxiliary obligations and the eight main cases $s=443,\ldots,450$. All tools read the LRAT stream directly, without an intervening proof conversion. The remaining eight cases $s=435,\ldots,442$ had two independently implemented checkers, without a formally verified checker. The complete theorem is computer-assisted; the set-to-formula reduction is not formalized.

The trusted base consists of the written cap induction, global symmetry and case-coverage arguments; the counting interpretation and its separately written audit program, which imports no encoder code; the proof checkers within the coverage stated; and the exhaustive witness checks. Each hypothetical set determines all auxiliary variables simultaneously by actual point counts. The auditor checks every clause under that one assignment rule, establishing that each set counterexample would extend to a satisfying assignment. In the replay procedure, the same audited CNF file is read by the solver and all checkers. Proof text is regenerated and checked through pipes. No solver verdict substitutes for a refutation check. The small-dimensional caps are proved within the computation, rather than imported as unchecked premises from Polymath.

Recent work addresses adjacent questions: Mouhib's symmetric and one-weight coloring bounds~\cite{mouhib}, Conlon--Farnsworth--Robertson's coloring questions~\cite{color}, Shelah's general bounds~\cite{shelah}, and Dahia's formalization of the general density theorem~\cite{dahia}. These are different endpoints and are not premises of this proof. This result determines one finite extremal value, not a new proof or improved general bound for the density theorem.

\section{Lines, caps and symmetry}
Any line-free set of at least 1303 points has a line-free 1303-point subset. Our formulas directly exclude the stronger condition of size at least 1303.

A wildcard pattern is a word in $\{0,1,2,*\}^n$ with at least one wildcard. Its line consists of $w(0),w(1),w(2)$, replacing all wildcards simultaneously. There are $4^n-3^n$ lines; thus the seven-dimensional audit checks 14,197 lines. A coordinate subcube is obtained by fixing some coordinates. A line-free set meets a $d$-dimensional subcube in at most $c_{d,3}$ points. Splitting on one coordinate gives
\begin{equation}\label{split} c_{d,3}\leq 3c_{d-1,3}.\end{equation}
For brevity put $c_d=c_{d,3}$ and $c_0=1$.

\begin{lemma}\label{sym}
A single permutation $\pi$ of the alphabet, applied simultaneously in every coordinate, preserves lines, line-freeness, cardinality and coordinate-subcube bounds. It permutes the three first-coordinate slice sizes arbitrarily.
\end{lemma}
\begin{proof}
Apply $\pi$ to each constant of a wildcard pattern, leaving its wildcard positions unchanged. The image of $w(a)$ is the resulting pattern evaluated at $\pi(a)$. Since $\pi$ is bijective, the image of the three-point line is exactly a line. The same map is bijective on all words and carries coordinate subcubes to subcubes of the same dimension. A first-coordinate slice labeled $a$ becomes the slice labeled $\pi(a)$. Choosing $\pi$ to send the three slice labels to their order by size proves the last assertion, including ties.
\end{proof}
Permuting symbols in the first coordinate alone is not the symmetry used here. For example, swapping 0 and 1 in that coordinate sends the line $\{00,11,22\}$ to $\{10,01,22\}$, which is not a combinatorial line. Earlier proof descriptions used this invalid local explanation. Lemma~\ref{sym} corrects it; the case formulas and the case cover are unchanged.

\section{The finite reduction}
The three auxiliary obligations are denoted $L_1,L_4,L_5$. They exclude line-free sets of sizes at least 3, 53 and 151 in dimensions 1, 4 and 5, respectively. Their formula inputs are:
\begin{center}
\begin{tabular}{llll}\toprule
Obligation & Dimension & Required size & Proper-subcube caps used\\\midrule
$L_1$ & 1 & 3 & none\\
$L_4$ & 4 & 53 & $2,6,18$\\
$L_5$ & 5 & 151 & $2,6,18,52$\\\bottomrule
\end{tabular}
\end{center}
The refutations establish $c_1\leq2$, then $c_4\leq52$, then $c_5\leq150$, in that order. Between them, (\ref{split}) establishes $c_2\leq6$, $c_3\leq18$ and $c_6\leq450$. There is no use of a cap of the dimension currently being proved. Thus no published small-dimensional optimality result is needed as an unchecked premise.

Suppose a line-free $A\subseteq\cube 7$ has $|A|\geq1303$. Write $a_v=|A\cap\{x_1=v\}|$. Lemma~\ref{sym} permits $a_0\geq a_1\geq a_2$. Since each slice is six-dimensional,
\[
435=\lceil1303/3\rceil\leq a_0\leq450.
\]
For each integer $s$ in this interval the main formula encodes line-freeness, size at least 1303, the proper-subcube caps $2,6,18,52,150,450$, and
\[
a_0=s\geq a_1\geq a_2.
\]
All sixteen formulas are unsatisfiable. Every counterexample to the upper bound would, after the global relabeling, satisfy one of them. Their refutations therefore imply $c_7\leq1302$.

\section{Encoding and independent semantic audit}
Point $x=(x_1,\ldots,x_n)$ has variable index $1+\sum_i x_i3^{n-i}$, interpreted as $x\in A$. Every line contributes the clause consisting of its three negated point variables. Hierarchical unary totalizers count points in coordinate subcubes and in intermediate unions of two child subcubes. Auxiliary variable $(S,j)$ means $|A\cap S|\geq j$, where $S$ is a product of nonempty subsets of the alphabet. Thresholds are truncated only beyond an available cap; the first forbidden threshold is explicitly false. A root threshold forces $|A|\geq K$. Units force $a_0=s$ and threshold implications force the slice order. The encoder and independent auditor are supplied as \texttt{checks/encode.py} and \texttt{checks/reduction.py}.

For clarity, the audit proves the direction required for an upper bound: every hypothetical counterexample extends to a satisfying assignment by the stated threshold interpretation. It does not rely on a solver's interpretation of auxiliary variables. For a box $S=M_1\times\cdots\times M_n$, let $d$ be its number of unrestricted coordinates. If $d<n$, disjoint coordinate subcubes give the valid bound
\[
U(S)=c_d\prod_{i:\,|M_i|<3}|M_i|.
\]
For the entire ambient cube use the trivial bound $3^n$, so the current claim is never its own premise. The empty set has bound zero.

The auditor recognizes every clause as one of the following valid statements under that interpretation.
\begin{enumerate}
\item A line clause forbids its three points together.
\item An upward counting clause has the form $(S,k)\vee\bigvee_r\neg(P_r,j_r)$, where the $P_r$ are disjoint subsets of $S$ and $\sum_r j_r\geq k$. Its validity follows by adding counts.
\item A downward counting clause has the form $\neg(S,k)\vee\bigvee_r(P_r,j_r)$. The remainder $S\setminus\bigcup_rP_r$ is empty or a box; disjointness and containment of the $P_r$ in $S$ are not required. By the union bound, when all positive literals are false,
\[
|A\cap S|\leq\sum_r\min(j_r-1,U(P_r))+U(S\setminus\textstyle\bigcup_rP_r)\leq k-1,
\]
which proves the clause.
\item A cap unit $\neg(S,k)$ has $U(S)\leq k-1$.
\item A premise clause follows from $|A|\geq K$, $a_0=s\geq a_1\geq a_2$: these include applicable positive root and slice thresholds, negative slice thresholds beyond $s$, and order implications $\neg(a_b\geq t)\vee(a_a\geq t')$ for $a<b$ and $t'\leq t$.
\end{enumerate}
The product sets and thresholds are read from an encoder-generated semantic map, but the auditor independently checks point indexing, the inclusion and disjointness conditions for upward clauses, bounds and every clause. The map supplies meanings, not evidence of their validity. The audit also checks all line clauses are present, the size and profile premises are present, header counts agree, and the exact integer case cover is complete. For the main formulas the rule counts are 14,197 line clauses, 940,716 upward clauses, 940,644 downward clauses, 765 cap clauses and 903 premise clauses (902 at $s=450$, where a cap already supplies one unit).

The logical dependency is therefore explicit: the written validity of these rules and the cap induction, the exhaustive clause audits, and checked unsatisfiability together rule out each set-theoretic counterexample. These steps have not been formalized together in a proof assistant.

\section{Refutations and the lower bound}
The exact tool revisions and build commands are in \texttt{checks/README.md}. CaDiCaL generates text LRAT proofs. The accepted computation regenerated all 19 refutations and checked every one with FRAT-rs's \texttt{lratchk}; cake\_lpr additionally checked the three auxiliary obligations and the eight main cases $s=443,\ldots,450$. The original computation used drat-trim's \texttt{lrat-check} on all 19, together with cake\_lpr on those 11 or a separate streaming checker on the other eight. This is not a claim that cake\_lpr checked all 19.

The eight obligations not replayed with cake\_lpr are the largest main cases, $s=435,\ldots,442$. Its heap usage grows with proof additions; the fixed-heap runs do not cover these proofs.

The public replay sends the generated stream to FRAT-rs and the supplied drat-trim checker, with cake\_lpr also enabled for those eleven obligations. An UNSAT exit code alone never passes: each applicable checker must accept an empty-clause derivation. The large proof streams are not distributed. Full regeneration produces about 1.56 TB of text, consumed through pipes; the accepted independent computation used about 64 solver CPU-hours, 9.7 FRAT-rs CPU-hours and 0.9 cake\_lpr CPU-hours, plus encoding and audits. A fresh public run also pays for drat-trim checking. The slowest case took about 21.4 hours elapsed. These are measured costs, not a promise for other hardware. The quick check audits the reduction and the witness and runs a small proof with corruption controls; it does not replay the sixteen main refutations.

For the lower bound, the supplied 1302 words (written over $\{1,2,3\}$) are a realization of Polymath's known lower bound~\cite[Section 2]{polymath}. They consist of all words with letter-count triples
\[
\begin{gathered}
034,043,124,142,223,232,250,304,\\
313,331,340,412,421,502.
\end{gathered}
\]
Their multinomial weights sum to 1302. Two independently written witness checkers verify distinctness, cardinality and absence of all 14,197 combinatorial lines. We do not identify this particular list with a specific source construction or claim a new lower bound. Together with the upper bound this proves Theorem~\ref{main}.

\paragraph{Availability and license.}
The accompanying bundle contains the manuscript source, the explicit witness, the formula specification, the encoder, semantic and symmetry audits, witness checkers, a standard LRAT checker and instructions for pinned external tools. Large proofs are regenerated. The work is licensed under Apache-2.0; the bundled drat-trim checker retains its MIT notice. See \texttt{THIRD\_PARTY.md} for dependencies and attribution.

\begin{thebibliography}{99}
\bibitem{fk} H. Furstenberg and Y. Katznelson, \emph{A density version of the Hales--Jewett theorem}, J. Anal. Math. 57 (1991), 64--119.
\bibitem{annals} D.H.J. Polymath, \emph{A new proof of the density Hales--Jewett theorem}, Ann. of Math. 175 (2012), 1283--1327, \href{https://doi.org/10.4007/annals.2012.175.3.6}{doi:10.4007/annals.2012.175.3.6}.
\bibitem{polymath} D.H.J. Polymath, \emph{Density Hales--Jewett and Moser numbers}, in An Irregular Mind, Bolyai Society Mathematical Studies 21, Springer, 2010, \href{https://arxiv.org/abs/1002.0374}{arXiv:1002.0374}; \href{https://doi.org/10.1007/978-3-642-14444-8_22}{doi:10.1007/978-3-642-14444-8\_22}.
\bibitem{tao2009} T. Tao, \emph{DHJ(3): 900--999 (Density Hales-Jewett type numbers)}, Polymath discussion post, March 4, 2009, \url{https://terrytao.wordpress.com/2009/03/04/dhj3-900-999-density-hales-jewett-type-numbers/}.
\bibitem{obryant} K. O'Bryant, \emph{Sets of Natural Numbers with Proscribed Subsets}, J. Integer Seq. 18 (2015), Article 15.7.7, \url{https://cs.uwaterloo.ca/journals/JIS/VOL18/OBryant/obr3.pdf}.
\bibitem{oeis} OEIS Foundation, \href{https://oeis.org/A156989}{A156989: density Hales--Jewett numbers for an alphabet of size 3}.
\bibitem{mouhib} Y. Mouhib, \emph{Lower Bounds for the Hales--Jewett Numbers via Symmetric and One-Weight Colorings}, \href{https://arxiv.org/abs/2606.22155v2}{arXiv:2606.22155v2}, 2026.
\bibitem{color} N. Conlon, N. Farnsworth and A. Robertson, \emph{On Hales--Jewett and Related Numbers}, \href{https://arxiv.org/abs/2607.18111}{arXiv:2607.18111}, 2026.
\bibitem{shelah} S. Shelah, \emph{HJ numbers revisited}, \href{https://arxiv.org/abs/2607.14732v3}{arXiv:2607.14732v3}, 2026.
\bibitem{dahia} G. Dahia, \emph{DensityHalesJewett}, formalization of the Dodos--Kanellopoulos--Tyros proof, \url{https://github.com/gdahia/DensityHalesJewett}, 2026.
\end{thebibliography}
\end{document}
